\documentclass{article}


% if you need to pass options to natbib, use, e.g.:
%     \PassOptionsToPackage{numbers, compress}{natbib}
% before loading neurips_2023


% ready for submission
\usepackage{neurips_2023}


% to compile a preprint version, e.g., for submission to arXiv, add add the
% [preprint] option:
%     \usepackage[preprint]{neurips_2023}


% to compile a camera-ready version, add the [final] option, e.g.:
%     \usepackage[final]{neurips_2023}


% to avoid loading the natbib package, add option nonatbib:
%    \usepackage[nonatbib]{neurips_2023}


\usepackage[utf8]{inputenc} % allow utf-8 input
\usepackage[T1]{fontenc}    % use 8-bit T1 fonts
\usepackage{hyperref}       % hyperlinks
\usepackage{url}            % simple URL typesetting
\usepackage{booktabs}       % professional-quality tables
\usepackage{amsfonts}       % blackboard math symbols
\usepackage{nicefrac}       % compact symbols for 1/2, etc.
\usepackage{microtype}      % microtypography
\usepackage{xcolor}         % colors


\title{COMP 560 Final Project: Spear Guy Survival}


% The \author macro works with any number of authors. There are two commands
% used to separate the names and addresses of multiple authors: \And and \AND.
%
% Using \And between authors leaves it to LaTeX to determine where to break the
% lines. Using \AND forces a line break at that point. So, if LaTeX puts 3 of 4
% authors names on the first line, and the last on the second line, try using
% \AND instead of \And before the third author name.

\usepackage{tikz,forest}
\usetikzlibrary{arrows.meta}

\forestset{
    .style={
        for tree={
            base=bottom,
            child anchor=north,
            align=center,
            s sep+=1cm,
    straight edge/.style={
        edge path={\noexpand\path[\forestoption{edge},thick,-{Latex}] 
        (!u.parent anchor) -- (.child anchor);}
    },
        }
    }
}


\begin{document}

\maketitle

\begin{abstract}
We present Spear Guy Survival, a 2D single-boss souls-like game developed in GameMaker Studio. The project focuses on fast-paced melee combat, player timing, and an adaptive boss encounter. The player uses multiple attacks, movement options, and blocking mechanics to survive against a spear-wielding enemy. To increase challenge over time, the boss uses Monte Carlo Tree Search (MCTS) to explore attack choices and reinforce successful action patterns. Our project demonstrates how lightweight AI methods can improve enemy behavior in a small-scale action game.
\end{abstract}\section{Introduction}

Designing a good game enemy is an important problem in artificial intelligence. In action games, especially boss fights, the enemy should feel challenging and interesting, but not unfair or too easy to predict. If a boss only follows a fixed pattern, players can quickly learn the same response every time, which makes the fight less exciting. Because of this, game AI is useful for creating enemies that can react to the player and make each encounter feel more dynamic.

In this project, we present \textit{Spear Guy Survival}, a 2D single-boss souls-like game developed in GameMaker Studio. The game focuses on fast combat, where the player must survive against a spear-wielding boss by using movement, attacks, and blocking at the right time. To make the encounter more engaging, we designed the boss to use Monte Carlo Tree Search-inspired decision making, so it can explore different attack choices and adapt its behavior based on what works during the fight. This project shows how AI methods can improve enemy behavior even in a small student game project.

\section{Game Design and Mechanics}
% Rian writes here\section{Boss AI with Monte Carlo Tree Search}

The boss AI in \textit{Spear Guy Survival} uses a simplified decision system inspired by Monte Carlo Tree Search. Instead of following a fixed attack loop, the boss chooses among four possible actions during combat. At the beginning, each action has the same initial weight, so the boss has no strong preference for one move over another. During gameplay, one action is selected at a time using a weighted random choice, which introduces variation while still allowing the system to favor actions that perform well.

After an action is chosen, the boss waits until the player is in the correct range and the attack cooldown has finished. Once the action is executed, the game checks whether it successfully damaged the player. If the move succeeds, its weight is increased; if it fails, its weight is reduced, while keeping a minimum weight so that all actions can still appear. In this way, the boss gradually changes its behavior based on the results of the fight and becomes more likely to reuse actions that are effective against the player.

Although this approach was motivated by ideas from Monte Carlo Tree Search, the final system is closer to an MCTS-inspired adaptive action policy than to a full classical MCTS implementation. A standard MCTS algorithm would build an explicit search tree, simulate future action sequences, and backpropagate rollout results through the tree. In our system, the adaptation happens online through immediate feedback from combat outcomes. This still captures the main idea of improving decisions through repeated outcomes, while keeping the boss responsive during the fight.

One important advantage of this system is that it makes the boss feel less scripted and more adaptive. Successful actions become more likely over time, so the encounter can change depending on how the player reacts, blocks, or positions during combat. At the same time, the weighted selection keeps some randomness in the boss behavior, which prevents the fight from becoming too repetitive. As a result, the boss can remain varied while also showing patterns that feel intentional and reactive.

\section{Implementation}
The enemy action is chosen by a function that takes in the weights for each action. All of the weights are added together to calculate the total possible actions, and the probability of each action is calculated by the weight of that action divided by that total sum. If a move succeeds, meaning damages the player, we add 1 to the action's weight. If it fails, meaning no damage done to the player, we subtract 1 from the value. Each action starts with a weight value of 2, which was chosen instead of 1 so that the initial state will change upon a failed action, and the minimum is 1 so that all actions remain possible. Below is an example of the decision tree the enemy creates upon the first action, where the first action is slash and it successfully damages the player. The percentages shown as edge labels are the percentage chance of choosing the action the edge leads to from the preceding state. Slash, thrust, upward thrust, counter, success, and failure, are abbreviated to Sl, T, Ut, C, S, and F, respectively.
% [V,edge label={node[midway,left,font=\scriptsize]{head}}]
\begin{center}
\begin{forest} 
[$Start$, tikz={\draw[{Latex}-, thick] (.north) --++ (0,1);}
    [$Sl$, edge label={node[midway,left,above=5pt,font=\scriptsize]{25\%}}
        [$S$
            [$Sl$, edge label={node[midway,left,above=5pt,font=\scriptsize]{33\%}}
            [$S$
            ]
            [$F$] 
        ]   
        [$T$, edge label={node[pos=0.8,left,font=\scriptsize]{22\%}}
            [$S$]
            [$F$] 
        ]
        [$Ut$, edge label={node[pos=0.8,left,font=\scriptsize]{22\%}}
            [$S$]
            [$F$] 
        ]
        [$C$, edge label={node[pos=0.8,right,above=1pt,font=\scriptsize]{22\%}}
            [$S$]
            [$F$] 
        ]
        ]
        [$F$] 
    ]   
    [$T$, edge label={node[pos=0.8,left,font=\scriptsize]{25\%}}
        [$S$]
        [$F$] 
    ]
    [$Ut$, edge label={node[pos=0.8,left,font=\scriptsize]{25\%}}
        [$S$]
        [$F$] 
    ]
    [$C$, edge label={node[midway,left,above=5pt,font=\scriptsize]{25\%}}
        [$S$]
        [$F$] 
    ]
] 
\end{forest}
\end{center}

\begin{center}
Figure 1: Enemy Decision Tree
\end{center}

\section{Results and Discussion}

Our project successfully created a playable 2D single-boss combat experience that emphasized timing, movement, and player decision-making. The combination of multiple player attacks, blocking mechanics, and boss pressure created encounters where careless actions were punished and patient play was rewarded. This helped achieve the intended souls-like style of gameplay, where players improve by learning patterns and adapting over repeated attempts.

The Monte Carlo Tree Search (MCTS) system also contributed to the difficulty of the encounter. Rather than relying only on fixed scripted behavior, the boss explored different attack options and reinforced successful action sequences. As a result, the boss became less predictable over time and forced the player to respond carefully.

From a development perspective, GameMaker Studio allowed rapid prototyping of combat mechanics, movement systems, and enemy behaviors. One challenge was balancing the boss so that it felt difficult without being unfair. Overall, the final system demonstrated a strong combination of gameplay mechanics and adaptive AI behavior.

\section{Conclusion}

In this paper, we presented \textit{Spear Guy Survival}, a 2D single-boss souls-like game developed in GameMaker Studio. The project combined responsive player combat mechanics with an adaptive boss powered by Monte Carlo Tree Search.

Our results show that lightweight AI methods such as MCTS can improve enemy encounters in small-scale games by increasing replayability and reducing predictability. Future work could expand the project through additional bosses, health systems, improved animations, and more sophisticated learning behavior.

\section*{References}

{
\small
Engine used: https://gamemaker.io/en \newline
Knight sprite: https://free-game-assets.itch.io/free-fantasy-knight \newline
Enemy sprite: https://raramenian.itch.io/the-impaler-7-attcks \newline
Background Sprite: https://saurabhkgp.itch.io/pixel-art-forest-background-simple-seamless-parallax-ready-for-2d-platformer-s




}

\end{document}
